\section*{Розділ VII. Взаємодія з іншими органами}
\addcontentsline{toc}{section}{Розділ VII. Взаємодія з іншими органами}

\subsection*{7.1. Взаємодія з Конференцією студентів Університету (КСУ)}
\addcontentsline{toc}{subsection}{7.1. Взаємодія з Конференцією студентів Університету (КСУ)}
    7.1.1. СР СМ є підзвітною та підконтрольною КСУ в межах, визначених Положенням про ОСС та цим Положенням.

    7.1.2. СР СМ щорічно подає детальний письмовий звіт про свою діяльність на розгляд КСУ. Звіт має містити інформацію про виконання плану роботи, реалізовані проєкти, фінансову діяльність (за наявності фінансування) та інші ключові аспекти роботи.

    7.1.3. КСУ має право вимагати від СР СМ надання позачергових звітів або інформації з окремих питань її діяльності. Такі запити КСУ є обов'язковими до виконання СР СМ у встановлені КСУ терміни.

\subsection*{7.2. Взаємодія зі Студентською уповноваженою делегацією (СУД)}
\addcontentsline{toc}{subsection}{7.2. Взаємодія зі Студентською уповноваженою делегацією (СУД)}
    7.2.1. СР СМ взаємодіє з СУД з метою забезпечення дотримання нормативних актів ОСС та захисту прав студентів.

    7.2.2. СР СМ зобов'язана надавати СУД на її запит будь-яку інформацію та документацію (протоколи, рішення, звіти, внутрішні документи тощо), необхідну для виконання СУД своїх повноважень, у терміни, встановлені СУД відповідно до Положення про СУД.

    7.2.3. Запити та рішення СУД, що потребують колегіального розгляду СР СМ, виносяться на розгляд найближчого засідання СР СМ. Голова СР СМ забезпечує підготовку відповіді або організацію виконання рішення СУД у термін, встановлений СУД.

    7.2.4. СР СМ бере до розгляду та виконання рекомендації та рішення СУД, прийняті в межах її компетенції.

\subsection*{7.3. Взаємодія з Центральною виборчою комісією студентів (ЦВКс)}
\addcontentsline{toc}{subsection}{7.3. Взаємодія з Центральною виборчою комісією студентів (ЦВКс)}
    7.3.1. СР СМ сприяє ЦВКс в організації та проведенні виборів до органів студентського самоврядування Університету.

    7.3.2. Сприяння полягає у:

        \begin{enumerate}[label=\alph*)]
            \item Наданні ЦВКс необхідної інформації для організації виборчого процесу (наприклад, контактні дані СРГ, інформація про структуру гуртожитків, дані про кількість мешканців тощо).
            \item Сприянні в інформаційному супроводі виборчої кампанії серед мешканців гуртожитків, поширенні оголошень та матеріалів ЦВКс через інформаційні ресурси СР СМ та СРГ.
            \item Наданні (за можливості та за попереднім погодженням) приміщень або технічних ресурсів, що знаходяться у розпорядженні СР СМ, для потреб ЦВКс під час виборчого процесу.
        \end{enumerate}

    7.3.3. СР СМ не втручається у здійснення ЦВКс її виключних повноважень щодо організації та проведення виборів.

\subsection*{7.4. Взаємодія зі Студентською радою Київського авіаційного інституту (СР КАІ)}
\addcontentsline{toc}{subsection}{7.4. Взаємодія зі Студентською радою Київського авіаційного інституту (СР КАІ)}
    7.4.1. СР СМ та СР КАІ взаємодіють на засадах партнерства та координації діяльності з метою представлення інтересів студентів, що мешкають у гуртожитках, та забезпечення належних умов проживання, навчання та дозвілля у студмістечку.

    7.4.2. Питання, що стосуються виключно внутрішнього розпорядку та умов проживання в гуртожитках, належать до компетенції СР СМ.

    7.4.3. Питання, що стосуються розвитку інфраструктури, безпеки, благоустрою території студмістечка або інші питання, що впливають на всіх студентів Університету, вирішуються спільно СР КАІ та СР СМ.

    7.4.4. У разі виникнення спорів щодо розмежування компетенції або неможливості досягнення згоди між СР СМ та СР КАІ, питання виноситься на розгляд КСУ.

\subsection*{7.5. Взаємодія з адміністрацією Університету та студмістечка}
\addcontentsline{toc}{subsection}{7.5. Взаємодія з адміністрацією Університету та студмістечка}
    7.5.1. Взаємодія СР СМ з адміністрацією Університету та студмістечка здійснюється відповідно до принципів та процедур, визначених Розділом VII Положення про ОСС, Статутом Університету та цим Положенням.

    7.5.2. Для забезпечення ефективного розгляду та підготовки проєктів рішень щодо погодження дій адміністрації з питань, зазначених у п. 7.2.5 Положення про ОСС та інших нормативних актах, в структурі СР СМ створюється постійно діюча \textbf{Комісія з питань поселення та взаємодії з адміністрацією студмістечка} (назва може бути уточнена окремим рішенням СР СМ).

    7.5.3. Порядок формування, склад та регламент роботи Комісії затверджуються СР СМ. Комісія здійснює попередній розгляд відповідних подань адміністрації, готує висновки та проєкти рішень для їх затвердження на засіданні СР СМ.

    7.5.4. Погодження рішень адміністрації з поточних (оперативних) питань здійснюється Головою СР СМ або Комісією за його дорученням. Погодження рішень адміністрації зі стратегічних питань (затвердження правил проживання, встановлення розміру плати за проживання, розподіл місць тощо) приймається колегіально на засіданні СР СМ. Перелік питань, що потребують колегіального погодження, затверджується рішенням СР СМ.

\subsection*{7.6. Взаємодія зі студентськими організаціями (СО)}
\addcontentsline{toc}{subsection}{7.6. Взаємодія зі студентськими організаціями (СО)}
    7.6.1. СР СМ підтримує діяльність студентських організацій (СО), зареєстрованих в Університеті у встановленому порядку, що здійснюють свою діяльність на території студмістечка, та може надавати їм інформаційну та організаційну підтримку.

\subsection*{7.7. Взаємодія з іншими органами та організаціями}
\addcontentsline{toc}{subsection}{7.7. Взаємодія з іншими органами та організаціями}
    7.7.1. СР СМ може взаємодіяти з комунальними службами, органами місцевого самоврядування, громадськими об'єднаннями та іншими організаціями з питань, що стосуються побуту, безпеки та інфраструктури студмістечка.
